\documentclass[10pt]{article}
\usepackage[margin=20mm]{geometry}
\usepackage{amsmath,amssymb,mathtools,fontspec,unicode-math,xurl,hyperref}
\IfFontExistsTF{DejaVu Serif}{\setmainfont{DejaVu Serif}}{\setmainfont{Latin Modern Roman}}
\setmathfont{latinmodern-math.otf}
\setlength{\parindent}{0pt}\setlength{\parskip}{6pt}\sloppy
\title{Scientific initial review of
The characteristic gluing problem for the Einstein equations and applications}
\author{M25-0145 Initial scientific review}
\date{}
\begin{document}
\maketitle
Abstract: A high-frequency, smooth-input sequence disproves the finite-Sobolev smoothness premise used in the nonlinear implicit-function argument. This is a Major issue in that proof route; it does not disprove existence of the asserted gluings or separately rule out a suitable gauge-slice argument. A second, local issue is the incorrect normalization of the displayed momentum-charge right inverse. Checks support the linear moment hierarchy, the correctly interpreted Kerr degree argument, the geometric decay recurrence and the causal angular-width calculation. The localization lemma’s support argument and several external analytic dependencies remain explicitly unverified.\par
\section*{Identity, version and scope}
The reviewed paper is S. Aretakis, S. Czimek and I. Rodnianski, The characteristic gluing problem for the Einstein equations and applications, arXiv:2107.02441v1, 6 July 2021, 31 pages [overview]. The assigned journal DOI is `10.1215/00127094-2024-0030`. The journal PDF was not obtained; no equivalence with its final text is asserted. All page references below are PDF pages, agreeing with printed pagination.\par
All 31 pages of the author manuscript were read as text. The bounded checks concerned the null-data definitions, linear obstruction count, gauge maps, nonlinear proof dependency, Kerr parameter matching, localized charge correction, radial iteration and geometric passage to spacelike data. The paper explicitly refers detailed linear and nonlinear estimates to its reference [10]. The corresponding arXiv:2107.02449v1, 102 pages [linear], was acquired and inspected only at selected definitions and proof steps needed for the gauge-map question. It was not subjected to a separate whole-paper review. Their actual roles and inspected material are recorded in the evidence bundle.\par
Theorem 1.1 concerns \(C^2\) gluing up to ten charges and gauge perturbations. Theorem 1.2 uses two null hypersurfaces and an auxiliary sphere for higher transverse regularity. These are different statements: the one-hypersurface result does not claim unrestricted higher-order transverse matching. Theorems 1.3–1.4 and Corollary 1.5 use Kerr parameter matching; Theorems 1.9–1.10 use localized correction and a rescaled iteration.\par
\section*{Established problems}
Problem 1 (Major: the full gauge-matching map is not Fréchet differentiable in the stated norms).  Section 3.2, p. 16 of [overview], asserts smooth dependence of sphere diffeomorphisms as a map \(\mathcal X(S)\times\mathcal Y_j\to\mathcal X(S)\). The sphere norm controls the lapse \(\Omega\) in \(H^6\), while \(\mathcal Y_j\) controls the diffeomorphism parameter in \(H^7\). The transversal parameter has the jet norm \[\|f\|_{\mathcal Y_f}=
 \|f|_S\|_{H^8}+\|f_u|_S\|_{H^6}
       +\|f_{uu}|_S\|_{H^4}+\|f_{uuu}|_S\|_{H^2}.\] The claimed smoothness fails, including for the composed gauge map with the original target fixed to exact Minkowski data.\par
A sequence within the prescribed parameter domain. On the round sphere, choose a real spherical harmonic \(A_n\) of degree and azimuthal order \(n\), normalized by \(\|A_n\|_{H^6}=1\). Its sine partner is orthogonal with the same norm. Put \[a_n=\frac{A_n}{n},\qquad
 f_n|_S=0,\quad (f_n)_u|_S=2a_n,\quad
 (f_n)_{uu}|_S=(f_n)_{uuu}|_S=0.\] These are smooth jets, realized, for example, by \(f_n(u,\theta)=2u
a_n(\theta)\) near \(u=0\), with a fixed smooth cutoff in \(u\). Thus \(\|f_n\|_{\mathcal Y_f}=2/n\). Let \(j_n\) represent the smooth azimuthal rotation \(R_{1/n}\); \(\|j_n\|_{\mathcal Y_j}=O(n^{-1})\). Sobolev embedding makes \(1+2a_n\) positive for large \(n\), so the coordinate change is admissible. The target characteristic data, the initial sphere data and the unknown unperturbed null data can all be held equal to the exact smooth Minkowski data.\par
The lapse transformation on p. 16, with the inverse-coordinate convention made explicit in [linear], p. 34, gives \[\Omega\bigl(P_{j_n}P_{f_n}(m)\bigr)
       =\sqrt{1+2a_n\circ R_{1/n}}.\] Reversing the convention for the angular pullback changes the sign of the rotation and leaves the following norm calculation unchanged. The derivative forced by the separate \(f\)- and \(j\)-directions has lapse variation \(a_n\); a rotation of the constant background lapse has zero variation. Since \(H^6(S^2)\) is a Banach algebra, \[\bigl\|\sqrt{1+2b}-1-b\bigr\|_{H^6}
       \leq C\|b\|_{H^6}^2\] for small \(b\). Rotation invariance and the azimuthal frequency give \[\|a_n\circ R_{1/n}-a_n\|_{H^6}
             =\frac{2\sin(1/2)}{n}.\] Consequently the remainder after the proposed linearization is at least \[\frac{2\sin(1/2)}{n}-O(n^{-2}),\] whereas the total parameter increment is \(O(n^{-1})\). Their ratio does not tend to zero. This contradicts Fréchet differentiability at the base point. Every member of the sequence is smooth; the obstruction is the specified finite-Sobolev topology.\par
The precise proof dependency. Section 3.1 of [overview], p. 15, reduces nonlinear gluing to an implicit-function theorem and points to [linear] for the detailed setup. Proposition 2.22 of that companion, p. 36, asserts the same smooth map. Its matching map \(\mathcal M\), Definition 2.11, p. 28, retains the entire lapse component. Definition 5.1, p. 63, contains \[F(x_{0,1},\widetilde x,x,f,j)
  =\bigl(x|_{S_{0,1}}-x_{0,1},\
       \mathcal M(x|_{S_{0,2}})-\mathcal M(P_jP_f(\widetilde x)),\
       (C_i(x))_i\bigr)\] and asserts full smoothness in an open Hilbert-space neighborhood. Holding the other variables at Minkowski leaves precisely the lapse remainder above. Thus the charge projection does not remove the counterexample, and giving the original target the stronger \(\mathcal X^+\) regularity does not resolve composition of the two free gauge parameters.\par
Counterevidence, impact and repair boundary. On pp. 64–65 the companion restricts \(F\) to the orthogonal complement of the kernel of its formal linearization. Rigid rotations are among the kernel directions. The sequence above has not been proved to lie in that complement, and does not independently disprove differentiability of the restricted map. The demonstrated defect is the false full-smoothness premise from which the text obtains the restricted smoothness needed for its implicit-function argument. This is directly tied to the nonlinear codimension-ten construction, hence Major under the contract’s core-derivation criterion.\par
A valid proof could establish the required regularity on a carefully chosen slice, change the gauge parametrization/order, or use a functional setup that handles derivative loss. Raising a norm alone also requires rechecking the bounded right inverse and all matching estimates. No complete such repair was established here. Neither Theorem 1.1 nor the subsequent existence theorems are declared false.\par
Problem 2 (Minor: the displayed localized momentum-charge prescription has the wrong normalization).  In Lemma 6.2, p. 23, equation (6.4) specifies a perturbation \(W\) with prescribed charges. Write \(r=1+2c\) and \(h=\Omega\operatorname{tr}\underline\chi\). For a target with \(\dot E_0=\dot L_0=\dot G_0=0\), \(\dot P_0^m=p\ne0\), the displayed formula sets \[\eta=0,\qquad h^{(1m)}=\frac{16\pi}{r^2\sqrt{4\pi}}p.\] The preceding charge formula instead gives \[\dot P^m(W)=\frac{\sqrt{4\pi/3}}{16\pi}r^2h^{(1m)}
                =\frac{p}{\sqrt3}\ne p.\]\par
This input is allowed: the required scalar moments can be realized by a smooth function supported in any fixed nonempty open transition region, as shown below. It is a counterexample to the explicit right-inverse prescription, not merely a normalization preference.\par
A correction, avoiding any ambiguity about vector-harmonic signs, is \[h^{(1m)}
  =\frac{16\pi\dot P_0^m/\sqrt{4\pi/3}
                +2r(\operatorname{div}\eta)^{(1m)}}{r^2},\] with the other moments chosen from the preceding charge identities. Support, regularity and radial powers in (6.1) are preserved. The correct charge map written immediately before (6.4) remains onto. Even the incorrectly labelled prescription has an invertible triangular effect on the charge parameters. Thus this local formula does not establish failure of the existence of a supported ten-charge adjustment or of Theorem 1.9. The supported impact is Minor.\par
\section*{Positive checks and rejected objections}
\subsection*{Linear charges and radial moments}
The gauge-invariant sector consists of the electric and magnetic degree-one vector modes of \(Q_0\), giving six components, and the degree-zero and degree-one scalar modes of \(Q_2\), giving four. This matches \(E,P,L,G\) on pp. 14, 16–17.\par
The shear moments used in (3.2)–(3.6) have distinct radial weights \(r^{-2},1,r^{-3},r^{-4}\). These functions are linearly independent on an open radial interval. Choosing a smooth positive interior weight \(w\), the moment Gram matrix \[G_{ij}=\int w(r)\,r^{p_i+p_j}\,\mathrm dr,
       \qquad (p_i)=( -2,0,-3,-4),\] is positive definite: its quadratic form is the integral of the square of a nonzero linear combination. This gives independent compactly supported controls for these finite moment conditions. The angular low-mode exceptions in the displayed elliptic operators agree with the separate conserved-charge treatment. This verifies the moment mechanism, not every inhomogeneous estimate or every higher-order commutator in the companion.\par
\subsection*{Kerr degree calculation}
A potential objection based on nonvanishing \(O(1)\) angular errors is incorrect. The source defines distinct \(O\) and little-\(o\) notation on p. 11; (4.7), p. 21, has \(o(1)\) in the last two components. For the parameter deviations define \[T_R(\delta E,P,\delta L,\delta C)
       =(\delta E,P,\delta L,\delta C-3RP).\] If the error in (4.7) is \(e=(O(R^{-1}),O(R^{-3/2}),o(1),o(1))\), then \[T_R^{-1}e=(e_E,e_P,e_L,e_G+3Re_P).\] After applying the ellipsoid weights \(\operatorname{diag}(R^{1/2},R^{1/2}I_3,I_3,I_3)\), all four components tend to zero, including the shear contribution \(3Re_P=O(R^{-1/2})\). Hence the degree argument is consistent, conditional on the stated uniform remainder estimates and on the positive-timelike Kerr parameter regime.\par
The strong spacelike decay \(k=O(r^{-5/2})\) implies zero limiting ADM momentum. Corollary 1.5 explicitly assumes positive ADM energy, so its limiting vector is timelike. Definitions (2.4)–(2.5) also allow \(M=0\), while the Kerr parameter construction in Section 4.2 assumes \(M>0\). The treatment of general zero-mass characteristic families is a remaining scope question; no impossibility example was proved. The detailed ADM/Kerr expansions and the spacelike existence theorem remain dependencies.\par
\subsection*{Supported moments and iteration}
For low scalar harmonics \(Y_i\) and a nonnegative smooth bump \(\chi\) supported in a fixed nonempty transition open set, the matrix \[B_{ij}=\int\chi\,Y_iY_j\,\mathrm d\mu\] is positive definite. A nonzero degree-at-most-one harmonic cannot vanish on an open set. Thus \(f=\chi\sum_i(B^{-1}v)_iY_i\) realizes arbitrary prescribed low moments. The same argument applies to the six analytic degree-one vector harmonics. Constants depend on the fixed open set; uniform bounds as its size shrinks are not asserted by this verification.\par
Put \(q=1+2c>1\). The exact charge recurrence in Section 7.2 has \[A_n=q^{-1}A_{n-1}+O(q^{-2n+2}R^{-2}).\] Multiplication by \(q^{n-1}\) makes the accumulated errors a convergent geometric sum. Therefore \(A_n=q^{-n+1}(A_1+O(R^{-2}))\), as actually parenthesized in (7.11). The \(L,G\) recurrence has coefficient \(q^{-2}\); its possible factor \(nq^{-2n}\) is bounded by a constant times \(q^{-n}\). There is no arbitrary \(C^n\) amplification in these exact recurrences. These are conditional checks of the iteration algebra, assuming the nonlinear local estimates and their uniformity.\par
For the causal calculation, a light-cone section is a ball of radius \(\rho c\), centered a distance \(\rho(1+c)\) from the origin. Its angular half-width is \(\arcsin(c/(1+c))\), independent of the annular scale \(\rho\). A fixed strict angular margin, followed by sufficiently large \(R\), is consistent with the domain-of-dependence construction. The characteristic well-posedness estimates were not independently reproved.\par
\section*{Other proof obligations and conventions}
Lemma 6.1, p. 23, claims that the solution orthogonal to the zero-endpoint kernel is supported in the angular support of its endpoint data. Its cutoff proof requires angular cutoff of free data to preserve zero endpoint data. The endpoint formulas contain angular differential operators, which generally have nonzero commutators with a cutoff. A justification specialized to this kernel and to the chosen orthogonal solution is still needed. No complete counterexample satisfying all bifurcate constraints and the orthogonality hypothesis was constructed. This remains Needs Evidence, not a proved false lemma.\par
The annular derivative norm in Section 7.3, p. 29, should be written using the current radius \(r_n=q^nR\): \[\sum_{\ell\leq m+1}r_n^\ell\|\partial^\ell(\widetilde g-e)\|_\infty
 +\sum_{\ell\leq m}r_n^{\ell+1}\|\partial^\ell\widetilde k\|_\infty
      \leq C r_n^{-1}.\] This implies the claimed \(r^{-1-\ell}\) and \(r^{-2-\ell}\) decay. The literal fixed-\(R\) weights printed there do not by themselves imply the second decay: \(k(r)=1/(Rr)\) meets that weaker annular estimate. This scalar example concerns the norm implication only, and is not vacuum initial data. The source’s preceding rescaling mechanism indicates the appropriate correction. Its first display with \(R\) in place of \(R^{-1}\) is merely weaker for \(R\geq1\), not an independent false inequality.\par
Two additional local conventions were checked to avoid building an argument on transcription errors. The geometric Schwarzschild expansions in the stated double-null convention are \(\Omega\operatorname{tr}\chi=2\Omega_M^2/r_M\) and \(\Omega\operatorname{tr}\underline\chi=-2\Omega_M^2/r_M\); the reference tuple on p. 13 prints different entries, including the wrong incoming sign already at \(M=0\). The later linearized formulas use the expected incoming sign. Also, the coordinate sentence \(u'=u+f\) on p. 15 is opposite to the pullback formulas on p. 16; the companion’s equation (2.46), p. 34, fixes the old coordinate as new coordinate plus \(f\). These require correction, but no separate false gluing theorem is inferred from them.\par
The extension, fill-in and exotic-data applications refer in part to a work listed as in preparation. Their full constructions, the complete higher-order bifurcate estimates, and the analytic details of the localized nonlinear function spaces were not verified in full here.\par
\section*{Conclusion and state}
Item | Initial-review status\par
Supported assessment | Major for the false full-smoothness premise of the nonlinear proof; Minor for the explicit charge normalization.\par
Repair verification | Charge normalization and scaling formulas checked; a complete nonlinear functional repair is not established.\par
Theorem validity | No counterexample to existence of the claimed gluings. A separate restricted-map proof remains possible.\par
Author adoption | Unknown; no author response requested.\par
Remaining work | Resolve the nonlinear functional setup, localization support lemma and zero-mass scope; compare final journal versions.\par
Scientific decision | Pending; human reviews: none.\par
The evidence bundle records exact input versions, selected companion pages, reviewer computations, auxiliary roles, actual compilation and rendered-page inspection. Source extraction and hash checks were not counted as proof verification. No author code was executed, and this report does not constitute independent acceptance.\par
2 S. Aretakis, S. Czimek and I. Rodnianski, The characteristic gluing problem for the Einstein equations and applications, arXiv:2107.02441v1 (6 July 2021), 31 pp. https://arxiv.org/abs/2107.02441v1. S. Aretakis, S. Czimek and I. Rodnianski, The characteristic gluing problem for the Einstein vacuum equations. Linear and non-linear analysis, arXiv:2107.02449v1 (6 July 2021), 102 pp. Selected definitions and proof dependencies only. https://arxiv.org/abs/2107.02449v1.\par
\end{document}
